\section{Introduction}
\label{sec:intro}

\paragraph{Motivation.}
Recent results on mathematical reasoning have shown that small open-source
language models, originally trained only with next-token prediction, can be
substantially improved by reinforcement learning with verifiable
rewards~\cite{deepseekmath, deepseekr1, dapo}. The recipe is now widely
agreed upon: (i) start from a competent base model, (ii) teach the model the
required answer format with a small supervised fine-tuning step, and
(iii) optimise a verifiable correctness signal on a much larger pool of
problems whose final answers can be checked automatically.

This report documents a course-project implementation of that recipe on a
single 1.5B-parameter base model, \modelname{Qwen2.5-1.5B}~\cite{qwen25}.
We deliberately keep the moving parts small: one base model, one SFT subset,
one main RL data pool with two scales, two RL algorithms (GRPO and DAPO),
and one shared evaluation suite. The goal is not to set a new
state-of-the-art number, but to (a) reproduce the reported pattern of
``small SFT followed by larger verifiable-reward RL'' end-to-end, (b) study
how each stage moves five evaluation benchmarks, and (c) check whether the
mathematical training causes any regression on general knowledge tasks.

\paragraph{Contributions.}
The concrete contributions of this report are:

\begin{enumerate}
  \item \textbf{End-to-end pipeline.} A reproducible three-stage pipeline
    (base $\rightarrow$ SFT $\rightarrow$ GRPO/DAPO) on
    \modelname{Qwen2.5-1.5B}, with strict separation between a clean SFT set
    requiring full reasoning traces and a large RL-ready set requiring only
    verifiable final answers (\cref{sec:data}).
  \item \textbf{Composite reward function.} A reward shape that combines
    exact-match and equivalence-aware correctness with several format,
    cleanliness, and length-control terms. The exact contributions are
    listed in \cref{sec:reward}.
  \item \textbf{Five-benchmark evaluation matrix.} A unified evaluation
    that covers three mathematical benchmarks (\code{GSM8K},
    \code{MATH-500}, \code{TheoremQA}) and two general benchmarks
    (\code{HellaSwag}, \code{MMLU}), all run through
    \emph{lm-eval-harness}~\cite{lmevalharness} with the same backend and
    decoding configuration. The complete numerical table is in
    \cref{sec:results}.
  \item \textbf{Failure-mode taxonomy.} A short qualitative analysis on
    saved \code{GSM8K} and \code{TheoremQA} samples that classifies the
    remaining errors into four recurring patterns
    (\cref{sec:failure}).
\end{enumerate}

\paragraph{Outline.}
\Cref{sec:data} describes data preparation and the four data layers.
\Cref{sec:method} explains the three-stage training design.
\Cref{sec:reward} details the composite reward.
\Cref{sec:exp} fixes the experimental setup and shows training-time
dynamics for the two RL variants. \Cref{sec:results} reports the main
benchmark table. \Cref{sec:general} focuses on whether mathematical
training degrades general ability. \Cref{sec:failure} analyses the
remaining failure modes, and \cref{sec:conclusion} summarises and lists
limitations.
